\section{Resultados}
\label{sec:resultados}

\subsection{Resposta individual de cada entrada}

A Tabela~\ref{tab:medidas} resume as seis simulações individuais,
comparando o ganho medido de cada entrada até a saída com o
coeficiente teórico da função global. Os ganhos parciais confirmam o
caminho de cada sinal: as entradas 1 a 3 aparecem somente em
$v_{out1}$ e as entradas 4 a 6 somente em $v_{out2}$. Além dos ganhos,
o nível DC medido na saída foi $-2{,}50$~V, em concordância com o valor
teórico $-12{,}5\times0{,}2 = -2{,}5$~V do termo constante do bloco A.

\begin{table}[htb]
  \centering
  \caption{Ganhos medidos por entrada (modelo ideal nível 1) e valores
  teóricos.}
  \label{tab:medidas}
  \begin{tabular}{lrrrr}
    \toprule
    Entrada & $v_{out1}/v_{in}$ & $v_{out2}/v_{in}$ &
    $v_{out}/v_{in}$ (sim.) & teórico \\
    \midrule
    $v_{in1}$ & $-9{,}00$ & $0$ & $+112{,}48$ & $+112{,}5$ \\
    $v_{in2}$ & $-3{,}00$ & $0$ & $+37{,}49$ & $+37{,}5$ \\
    $v_{in3}$ & $+2{,}50$ & $0$ & $-31{,}24$ & $-31{,}25$ \\
    $v_{in4}$ & $0$ & $+2{,}00$ & $+25{,}00$ & $+25$ \\
    $v_{in5}$ & $0$ & $+4{,}00$ & $+49{,}99$ & $+50$ \\
    $v_{in6}$ & $0$ & $-9{,}00$ & $-112{,}48$ & $-112{,}5$ \\
    \bottomrule
  \end{tabular}
\end{table}

A Figura~\ref{fig:individual} apresenta os dois casos de maior ganho:
$v_{in1}$ (ganho $+112{,}5$) e $v_{in6}$ (ganho $-112{,}5$). A
amplitude de \SI{5.625}{\volt}, as fases opostas e o nível DC de
$-2{,}5$~V (as senoides oscilam em torno desse nível, não do zero)
estão de acordo com os valores previstos.

\begin{figure}[htb]
  \centering
  \includegraphics[width=0.49\textwidth]{rs2_vin1_out.png}
  \includegraphics[width=0.49\textwidth]{rs2_vin6_out.png}
  \caption{Saída com somente $v_{in1}$ ativa (esquerda) e somente
  $v_{in6}$ ativa (direita), seno de \SI{50}{\milli\volt} e
  \SI{1}{\kilo\hertz}.}
  \label{fig:individual}
\end{figure}

A Figura~\ref{fig:individual-resto} apresenta as outras quatro
respostas individuais. Em cada uma, a amplitude da saída dividida
pelos \SI{50}{\milli\volt} da entrada corresponde ao ganho medido da
Tabela~\ref{tab:medidas}, e a fase em relação ao seno de entrada
indica o sinal do coeficiente: $v_{in2}$, $v_{in4}$ e $v_{in5}$
aparecem em fase (coeficiente positivo) e $v_{in3}$ em oposição de
fase (coeficiente negativo).

\begin{figure}[htb]
  \centering
  \includegraphics[width=0.49\textwidth]{rs2_vin2_out.png}
  \includegraphics[width=0.49\textwidth]{rs2_vin3_out.png}\\
  \includegraphics[width=0.49\textwidth]{rs2_vin4_out.png}
  \includegraphics[width=0.49\textwidth]{rs2_vin5_out.png}
  \caption{Saída com somente $v_{in2}$ (acima, à esquerda), $v_{in3}$
  (acima, à direita), $v_{in4}$ (abaixo, à esquerda) e $v_{in5}$
  (abaixo, à direita) ativas.}
  \label{fig:individual-resto}
\end{figure}

\subsection{Validação dos blocos}
\label{sec:res-blocos}

Nos ensaios por bloco, a saída de cada bloco foi comparada ponto a
ponto com a respectiva função teórica, calculada em paralelo por uma
fonte comportamental. O erro máximo foi $2{,}07\times10^{-4}$~V no
bloco A e $1{,}59\times10^{-4}$~V no bloco B, compatível com o erro
residual esperado para um ganho de malha aberta finito de
$2\times10^5$. No bloco C o ganho diferencial medido foi $+12{,}50$
(teórico $12{,}5$) e, com \SI{1}{\volt} de modo comum aplicado às duas
entradas, a média da saída ficou em $2{,}81\times10^{-8}$~V,
confirmando a rejeição de modo comum esperada para resistores
perfeitamente casados na simulação. A Figura~\ref{fig:blocoab}
apresenta as formas de onda dos blocos A e B.

\begin{figure}[htb]
  \centering
  \includegraphics[width=0.49\textwidth]{rs2_bloco_a_vout1.png}
  \includegraphics[width=0.49\textwidth]{rs2_bloco_b_vout2.png}
  \caption{Saída do bloco A (esquerda) e do bloco B (direita) no teste
  com vários senos somados.}
  \label{fig:blocoab}
\end{figure}

\subsection{Função global e macromodelo comercial}
\label{sec:res-global}

Com as seis entradas ativas simultaneamente, a saída acompanhou a
referência ideal (que inclui o nível DC de $-2{,}5$~V) com erro máximo
de $9{,}27\times10^{-4}$~V e erro RMS de $3{,}94\times10^{-4}$~V, para
uma excursão de $-8{,}20$ a $+3{,}20$~V
(Figura~\ref{fig:global}). Repetindo o ensaio com o macromodelo
comercial do LM741, o erro máximo subiu para $0{,}27$~V (RMS de
$0{,}13$~V), cerca de 2{,}4\% da excursão total de \SI{11.4}{\volt}. A
causa dominante é o polo do LM741 real em \SI{5}{\hertz}: nas
frequências do ensaio o ganho de malha aberta cai para a ordem de
$10^3$, aumentando o erro de cada estágio em relação ao modelo ideal.
Soma-se a esse efeito o offset de entrada de \SI{1}{\milli\volt},
multiplicado pelos ganhos até a saída. O circuito mantém o
comportamento projetado, como mostra a Figura~\ref{fig:vendor}.

\begin{figure}[htb]
  \centering
  \includegraphics[width=\textwidth]{rs2_completo_overlay.png}
  \caption{Ensaio global: saída do circuito sobreposta à referência
  ideal.}
  \label{fig:global}
\end{figure}

\begin{figure}[htb]
  \centering
  \includegraphics[width=\textwidth]{rs2_vendor_overlay.png}
  \caption{Ensaio global com o macromodelo comercial do LM741.}
  \label{fig:vendor}
\end{figure}

Como verificação adicional, o mesmo ensaio global foi executado com um
amplificador operacional moderno, o ADA4610 (entrada JFET, modelo
embarcado no LTspice), com a mesma alimentação de
$\pm$\SI{15}{\volt}. O erro máximo caiu de $0{,}27$~V para
$8{,}7\times10^{-3}$~V (RMS de $3{,}4\times10^{-3}$~V), cerca de 30
vezes menor que o obtido com o LM741 e inferior a 0{,}1\% da excursão
(Figura~\ref{fig:moderno}). Esse resultado confirma que, mantido o ganho de malha aberta elevado
nas frequências do ensaio, o erro residual por estágio retorna ao
nível obtido com o modelo ideal. Para este circuito, a substituição
do LM741 por um amplificador moderno reduz o desvio observado no
ensaio anterior a níveis desprezíveis.

\begin{figure}[htb]
  \centering
  \includegraphics[width=\textwidth]{rs2_moderno_overlay.png}
  \caption{Ensaio global com o amplificador moderno ADA4610.}
  \label{fig:moderno}
\end{figure}

\subsection{Discussão}
\label{sec:discussao}

A cascata de somadores inversores atende à restrição do bloco A sem
divisores na entrada não inversora: o sinal positivo de $v_{in3}$
resulta da dupla inversão e cada ganho depende de uma única razão de
resistores. O amplificador de instrumentação conecta os dois blocos
sem carregá-los, em razão da alta impedância do seu estágio de
entrada.

Duas limitações práticas devem ser registradas. Os resistores de
$100/9$ e $100/3$~\si{\kilo\ohm} não pertencem à série E96; em uma
montagem física seriam necessárias associações (por exemplo, 11k em
série com 110~$\Omega$, e 33k em série com 330~$\Omega$), com erro da
ordem da tolerância dos componentes. Além disso, os ganhos elevados
(módulo $112{,}5$) limitam a amplitude de entrada: acima de cerca de
\SI{100}{\milli\volt} de pico em $v_{in1}$ ou $v_{in6}$ a saída
satura no semiciclo negativo, pois o nível DC de $-2{,}5$~V reduz a
margem nesse sentido. Todos os ensaios respeitaram esse limite.
